\documentclass[12pt,a4paper]{article}
\usepackage[utf8]{inputenc}
\usepackage[T1]{fontenc}
\usepackage{lmodern}
\IfFileExists{brazil.ldf}{\usepackage[brazil]{babel}}{%
  \usepackage{babel}\babelprovide[main,import]{portuguese}}
\usepackage[margin=2.5cm,headheight=16.6pt]{geometry}
\usepackage{array,tabularx,booktabs}
\usepackage{amsmath}
\usepackage{enumitem}
\usepackage{fancyhdr}
\usepackage{xcolor}
\usepackage{titlesec}
\definecolor{azulProfessor}{HTML}{184C78}
\definecolor{cinzaTitulos}{HTML}{1E1E1E}
\definecolor{cinzaCabecalho}{gray}{0.5}
\titleformat{\section}{\Large\bfseries\color{azulProfessor}}{\thesection.}{0.6em}{}[\titlerule]
\titleformat{\subsection}{\large\bfseries\raggedright\color{cinzaTitulos}}{\thesubsection.}{0.6em}{}
\titleformat{\subsubsection}{\normalsize\bfseries\raggedright\color{cinzaTitulos}}{\thesubsubsection.}{0.6em}{}
\titlespacing*{\section}{0pt}{3.5ex plus 1ex minus .2ex}{2.3ex plus .2ex}
\titlespacing*{\subsection}{0pt}{3.25ex plus 1ex minus .2ex}{1.5ex plus .2ex}
\titlespacing*{\subsubsection}{0pt}{2.5ex plus 1ex minus .2ex}{1.2ex plus .2ex}
\setcounter{tocdepth}{2}
\usepackage{hyperref}
\hypersetup{colorlinks=true,linkcolor=azulProfessor,urlcolor=azulProfessor,citecolor=azulProfessor,linktoc=section,pdftitle={Laboratório Educativo de IA Generativa com GANs — Relatório Final — Sprint 2},pdfauthor={Allan Vitor Piran Lemes e Guilherme Rodrigues dos Santos}}
\urlstyle{same}
\newcolumntype{Y}{>{\raggedright\arraybackslash}X}
\setlength{\parindent}{1.5em}
\setlength{\parskip}{0pt}
\setlength{\emergencystretch}{2em}
\renewcommand{\arraystretch}{1.2}
\setlist[itemize]{leftmargin=*,itemsep=5pt,topsep=6pt}
\pagestyle{fancy}
\fancyhf{}
\fancyhead[L]{\small\color{cinzaCabecalho}Oficina de Projetos em Ciência de Dados e IA}
\fancyhead[R]{\small\color{cinzaCabecalho}Relatório Final -- UTFPR}
\fancyfoot[L]{\small\color{cinzaCabecalho}Documento de Ciência de Dados \& IA}
\fancyfoot[R]{\small Página \thepage}
\renewcommand{\headrulewidth}{0.4pt}
\renewcommand{\footrulewidth}{0.4pt}
\newcommand{\pendente}[1]{\textit{#1}}

\begin{document}
\renewcommand{\contentsname}{Sumário}
\renewcommand{\refname}{Referências}

\begin{titlepage}
\thispagestyle{fancy}
\centering
\vspace*{0.7cm}
{\Large\bfseries UNIVERSIDADE TECNOLÓGICA FEDERAL DO PARANÁ\par}
\vspace{0.15cm}
{\large DEPARTAMENTO ACADÊMICO DE INFORMÁTICA\par}
\vspace{0.1cm}
{\small CAMPUS PATO BRANCO\par}
\vspace{0.1cm}
{\small OFICINA DE PROJETOS EM CIÊNCIA DE DADOS E IA\par}
\vspace{1.6cm}
\noindent\rule{\textwidth}{1pt}\par
\vspace{0.2cm}
{\Huge\bfseries\color{azulProfessor}Laboratório Educativo de IA\par Generativa com GANs\par}
\vspace{0.35cm}
{\large\bfseries RELATÓRIO FINAL DE PROJETO\par}
\vspace{0.25cm}
{\small Atualização parcial: encerramento da Sprint 2\par}
\vspace{0.45cm}
\noindent\rule{\textwidth}{1pt}\par
\vspace{1.2cm}
{\renewcommand{\arraystretch}{1.05}
\begin{tabular}{@{}ll@{}}
\textbf{Equipe de Desenvolvimento:} & Allan Vitor Piran Lemes \\
 & Guilherme Rodrigues dos Santos \\
\textbf{Orientador / Prof.:} & Dr. Luiz Fernando Puttow Southier \\
\textbf{Data de Atualização:} & 30 de outubro de 2026 \\
\textbf{Versão do Documento:} & 2.0 \\
\end{tabular}\par}
\end{titlepage}
\setcounter{page}{2}
\tableofcontents
\clearpage

\section{Introdução}
Este relatório acompanha o desenvolvimento do Laboratório Educativo de IA Generativa com GANs, proposto para apresentar conceitos de inteligência artificial a estudantes do ensino médio por meio de uma apostila e de uma aplicação local. O projeto é desenvolvido por dois estudantes, com atividades de estudo, implementação e documentação distribuídas nas sprints previstas no Plano de Projeto entregue.

A proposta consiste em utilizar uma Rede Adversarial Generativa (GAN) treinada com o Fashion-MNIST, uma base de imagens de roupas, calçados e acessórios. As imagens possuem resolução de $28\times28$ pixels, em tons de cinza~\cite{fashionmnist}. A aplicação educativa deverá mostrar exemplos da base, permitir a comparação de imagens geradas em diferentes momentos do treinamento e oferecer geração, visualização e download de imagens.

A utilidade do produto está em tornar o funcionamento da IA generativa mais acessível por meio de uma demonstração prática. A baixa resolução será apresentada como característica e limitação do modelo. O roteiro previsto terá duração inicial de 20 a 30 minutos e orientará a leitura e exploração autônoma por leitores iniciantes. O treinamento será realizado previamente pela equipe.

\textbf{Esta versão reúne os registros das Sprints 1 e 2, com atualização em 30/10/2026.} A Sprint 1, de 02/10 a 09/10, concentrou-se na organização do trabalho e na preparação dos dados. Na Sprint 2, de 09/10 a 30/10, foram implementados o treinamento da GAN, o salvamento do progresso e a geração de imagens a partir do modelo salvo. O registro técnico da etapa apresenta uma implementação inicial do Allan e uma adaptação local posterior, cujos resultados são tratados separadamente~\cite{registroS2}. A interface e a apostila permanecem para as próximas etapas.


\section{Objetivos do Projeto}
\subsection{Objetivos}
O objetivo geral é desenvolver um laboratório educativo com uma GAN treinada no Fashion-MNIST, para apoiar a introdução à IA generativa no ensino médio por meio de materiais introdutórios, registrando o desenvolvimento e os aprendizados da equipe.

Os objetivos específicos do projeto são:
\begin{itemize}
\item Estudar o funcionamento da GAN, identificando o papel do gerador e do discriminador, e conhecer a base de dados utilizada.
\item Preparar o ambiente e adaptar uma implementação introdutória ao projeto.
\item Treinar o modelo e registrar a evolução visual das imagens geradas.
\item Construir uma interface simples para conhecer a base, comparar momentos do treinamento e gerar, visualizar e baixar imagens.
\item Documentar o código, as atividades, as dificuldades e os aprendizados de cada sprint.
\item Preparar material pedagógico e um roteiro de uso voltados ao ensino médio e a leitores iniciantes.
\end{itemize}

Esses objetivos abrangem o projeto completo. A Sprint 1 estabeleceu a organização e a preparação dos dados; a Sprint 2 avançou na implementação, no primeiro treinamento e na geração de imagens com o modelo salvo. A versão atual recebe números aleatórios, sem descrição textual da roupa ou seleção de categoria. A interface educativa e a seleção de categorias ainda não foram implementadas.

\subsection{Restrições Técnicas, Regulatórias e Éticas}
O desenvolvimento deve ser compatível com o tempo da disciplina, os recursos disponíveis e os conhecimentos adquiridos pelos integrantes durante o trabalho. Serão utilizados computadores pessoais e recursos gratuitos das ferramentas escolhidas. O uso de um ambiente de treinamento em nuvem será avaliado apenas se necessário e conforme sua disponibilidade.

O planejamento considera dedicação média de cinco horas semanais por integrante. As quatro sprints abrangem nove semanas, com mais duas semanas previstas para a consolidação do relatório e da apresentação. A data final de 18/12/2026 permanece sujeita à confirmação.

As fontes de dados e os exemplos de código efetivamente utilizados deverão ser identificados. As imagens geradas serão apresentadas como resultados de IA, incluindo suas limitações. A avaliação prevista é descritiva: funcionamento da aplicação, evolução visual e revisão dos materiais pelos autores. Não foi estabelecida uma porcentagem de acerto para a GAN.

As 10.000 imagens do conjunto oficial de teste do Fashion-MNIST permaneceram fora do treinamento nas duas versões descritas. A divisão adicional de 50.000 imagens para treinamento e 10.000 para validação, preparada na Sprint 1, passou a ser utilizada efetivamente no treinamento local da Sprint 2, conforme a Seção~\ref{sec:ajustesS2}. Não foi implementada uma métrica quantitativa de validação nesta etapa.

\subsection{Avaliação e Acurácia das Estimativas (Planejado vs. Executado)}
\label{sec:estimativas}
O plano adotou uma estimativa baseada na capacidade dos dois integrantes, com dedicação média de cinco horas semanais por pessoa. Foram previstas 110 horas-pessoa para o projeto: 90 horas nas quatro sprints e 20 horas para a etapa final. A Sprint 1 possui previsão de 10 horas-pessoa, incluindo o planejamento inicial; a Sprint 2 possui previsão de 30 horas-pessoa.

Na Sprint 1, foram informadas 3 horas de trabalho do Allan e 2 horas do Guilherme, totalizando 5 horas. As horas efetivamente dedicadas à Sprint 2 ainda não foram consolidadas. A tabela distingue a capacidade prevista do esforço informado.

\begin{center}
\small
\begin{tabularx}{\textwidth}{@{}YrrY@{}}
\toprule
\textbf{Etapa} & \textbf{Previsto (h)} & \textbf{Realizado (h)} & \textbf{Desvio} \\
\midrule
Sprint 1 & 10 & 5 & $-50\%$ \\
Sprint 2 & 30 & --- & Não apurado \\
\midrule
Acumulado até a Sprint 2 & 40 & --- & Não apurado \\
\bottomrule
\end{tabularx}
\end{center}

A variação é calculada por:
\[
\text{Desvio (\%)} =
\frac{\text{Horas realizadas} - \text{Horas previstas}}
     {\text{Horas previstas}} \times 100.
\]
Assim, o desvio da Sprint 1 é $(5-10)/10\times100=-50\%$. Isso indica esforço registrado inferior à capacidade prevista, sem permitir concluir, por si só, que todas as tarefas foram atendidas ou que houve ganho de produtividade. Não foi informado um motivo específico para a diferença. O desvio da Sprint 2 e o acumulado dependem do registro de suas horas.

O valor ilustrativo de R\$ 20,00 por hora foi mantido. O custo previsto é de R\$ 200,00 para a Sprint 1 e R\$ 600,00 para a Sprint 2, totalizando R\$ 800,00 previstos até esta etapa. As 5 horas informadas da Sprint 1 correspondem a R\$ 100,00 nesse exercício de cálculo. O equivalente realizado da Sprint 2 ainda não pode ser apurado. A previsão global permanece em R\$ 2.200,00.

Esses valores não representam pagamento aos estudantes. O desembolso financeiro previsto continua em R\$ 0,00, com uso dos recursos disponíveis; não foi apresentado um levantamento de despesas realizadas. Os 296,8 segundos medidos no treinamento local representam tempo de máquina e não substituem o registro de dedicação ativa dos integrantes.

\section{Gestão de Riscos do Projeto}
A análise dos quatro riscos do planejamento foi atualizada até o encerramento da Sprint 2. As situações abaixo se baseiam nas atividades e nos resultados registrados, sem antecipar uma avaliação final do projeto.

\subsection*{R01 — Dificuldade para compreender ou executar a GAN}
\textbf{Previsão:} probabilidade alta e impacto alto. A prevenção consiste em estudar um exemplo introdutório em conjunto e testar cada etapa. A contingência é buscar apoio do professor e adotar uma implementação menor, preservando a geração básica.

\textbf{Situação até a Sprint 2:} a implementação executou um treinamento completo de uma época e gerou imagens com o modelo salvo. Foi necessário ajustar o uso dos dados para que o treinamento respeitasse a divisão preparada na Sprint 1. Os testes curtos ajudaram a verificar execução, salvamento e retomada. O registro não detalha dificuldades individuais de compreensão; portanto, não é possível avaliar esse aspecto separadamente.

\subsection*{R02 — Treinamento demorado ou interrompido}
\textbf{Previsão:} probabilidade média e impacto médio. Estão previstos um teste curto antes de execuções maiores e o salvamento do progresso. A contingência é retomar o modelo salvo ou avaliar outro ambiente gratuito disponível.

\textbf{Situação até a Sprint 2:} a execução local de uma época terminou em 296,8 segundos, em CPU, sem interrupção registrada. O salvamento e a retomada foram verificados em testes curtos. Essa verificação não demonstra a duração ou a estabilidade de um treinamento completo de 15 épocas. A necessidade de outro ambiente ainda não foi estabelecida.

\subsection*{R03 — Imagens pouco reconhecíveis}
\textbf{Previsão:} probabilidade alta e impacto médio. A prevenção é conferir os dados e acompanhar as imagens desde os primeiros testes. A contingência é revisar a implementação, realizar ajustes simples e apresentar as limitações observadas.

\textbf{Situação até a Sprint 2:} o risco se manifestou na análise visual relatada. Após uma época, surgiram contornos que lembram roupas e calçados, mas persistiram borramento, padrões quadriculados e amostras ambíguas. A visualização foi padronizada para permitir comparações entre épocas da execução local. As imagens ainda apresentam limitações de nitidez.

\subsection*{R04 — Falta de tempo para concluir as atividades}
\textbf{Previsão:} probabilidade média e impacto alto. A prevenção é acompanhar as tarefas no Jira e comparar o trabalho restante com a capacidade da sprint. A contingência é priorizar as funções básicas, a demonstração educativa e a documentação.

\textbf{Situação até a Sprint 2:} a equipe informou a inclusão das tarefas no backlog do Jira e o encerramento da sprint. A entrega técnica inicial foi alcançada. Permanecem sem registro as horas desta etapa. Não há informação suficiente para atribuir essas pendências à falta de tempo; elas deverão ser acompanhadas na continuidade do projeto.

\section{Execução por Sprints e Alocação de Recursos}
O desenvolvimento está organizado em quatro sprints e uma etapa de consolidação final. A equipe é formada por Allan Vitor Piran Lemes e Guilherme Rodrigues dos Santos. Os percentuais associados às entregas representam pesos de avaliação da disciplina, e não percentuais de conclusão ou de esforço.

\subsection{Primeira Sprint: Organização, Jira e Contato com os Dados}
\textbf{Período:} 02/10 a 09/10/2026. \textbf{Peso da entrega:} 3\%.

O objetivo desta sprint foi estabelecer a organização inicial do trabalho, preparar o acesso à base Fashion-MNIST e relacionar a proposta técnica à linguagem acessível ao ensino médio. As atividades descritas a seguir correspondem ao registro de execução da equipe.

\subsubsection{Atividades realizadas}
Foi configurado o repositório Git no GitHub, incluindo o arquivo \path{.gitignore}, para organizar o código e controlar quais arquivos serão acompanhados no versionamento. Também foi estruturado o quadro inicial do Jira, com as quatro sprints e os respectivos prazos. Essa organização servirá de referência para distribuir e acompanhar as próximas tarefas.

Na preparação dos dados, foi desenvolvido o script Python \path{preparacao_dados.py} para baixar o Fashion-MNIST. O registro da sprint informa a separação de imagens para treinamento e validação e a exportação de um painel inicial de amostras da base no diretório \path{artefatos/}. As características dessa organização são apresentadas a seguir.

Na dimensão educativa, foi redigido o objetivo do material introdutório com linguagem acessível ao ensino médio.

\subsubsection{Organização inicial dos dados}
\label{sec:dados}
O Fashion-MNIST possui 60.000 imagens no conjunto oficial de treinamento e 10.000 no conjunto oficial de teste, distribuídas em dez categorias~\cite{fashionmnist}. Nesta sprint, o conjunto de treinamento foi subdividido em 50.000 imagens para treinamento e 10.000 para validação, com divisão estratificada e semente 42. A estratificação busca manter a proporção das categorias nos dois subconjuntos; a semente permite repetir a mesma divisão.

\begin{center}
\small
\begin{tabularx}{\textwidth}{@{}YrY@{}}
\toprule
\textbf{Conjunto} & \textbf{Quantidade} & \textbf{Origem e destinação} \\
\midrule
Treinamento & 50.000 & Parte do treinamento oficial, destinada ao aprendizado da GAN. \\
Validação & 10.000 & Parte reservada do treinamento oficial para verificações posteriores. \\
Teste oficial & 10.000 & Conjunto original separado, mantido fora do treinamento. \\
\bottomrule
\end{tabularx}
\end{center}

Essa subdivisão é um detalhamento da preparação dos dados que não estava especificado no plano. Ela mantém a separação do teste oficial. Ao final da Sprint 1, essa divisão ainda não constituía um resultado de treinamento ou de avaliação da GAN. Sua integração ao treinamento é relatada na Sprint 2. O painel exportado contém exemplos da base, e não imagens produzidas pela GAN.

\subsubsection{Comparação entre o planejado e o realizado}
\label{sec:comparacao}
\begin{center}
\small
\begin{tabularx}{\textwidth}{@{}>{\raggedright\arraybackslash}p{0.29\textwidth}Y@{}}
\toprule
\textbf{Previsão da Sprint 1} & \textbf{Registro no encerramento} \\
\midrule
Organizar o repositório e preparar o ambiente & Repositório configurado, com \path{.gitignore}, e preparação dos dados realizada em Python nos computadores da equipe. \\
Organizar o Jira e registrar os prazos & Quadro inicial configurado com quatro sprints e prazos. O registro disponível não detalha os cartões e seus critérios de conclusão. \\
Definir tarefas e responsáveis & Divisão principal entre os integrantes registrada. O detalhamento dos cartões deverá ser vinculado como evidência. \\
Baixar a base e visualizar exemplos & Script de preparação desenvolvido, divisão dos dados realizada e painel de amostras exportado. \\
Descrever a proposta educativa & Objetivo educativo e leitores pretendidos definidos. \\
Registrar atividades, dúvidas e aprendizados & Atividades e contribuições registradas, com 3 horas do Allan e 2 horas do Guilherme. Faltam as reflexões individuais e a inclusão das evidências. \\
\bottomrule
\end{tabularx}
\end{center}

\subsubsection{Recursos utilizados e divisão das tarefas}
Foram utilizados os computadores pessoais dos integrantes. As ferramentas informadas nesta sprint foram Python 3.10 ou superior, PyTorch, Torchvision, Scikit-learn, Git/GitHub, Jira e Overleaf. Não há uso de infraestrutura de treinamento em nuvem registrado nesta etapa.

\begin{center}
\small
\begin{tabularx}{\textwidth}{@{}>{\raggedright\arraybackslash}p{0.28\textwidth}Y@{}}
\toprule
\textbf{Integrante} & \textbf{Contribuição na Sprint 1} \\
\midrule
Allan Vitor Piran Lemes & Organização do repositório, desenvolvimento da preparação dos dados e caracterização dos leitores da apostila. \\
Guilherme Rodrigues dos Santos & Estruturação do quadro inicial do Jira e planejamento das tarefas das sprints. \\
\bottomrule
\end{tabularx}
\end{center}

A distribuição das atividades está de acordo com a divisão principal prevista para a primeira sprint no planejamento entregue. A organização adotada permite que cada integrante conduza suas tarefas e acompanhe o trabalho do colega. As horas informadas foram consolidadas na Seção~\ref{sec:estimativas}.

\subsubsection{Aprendizados, dificuldades e pendências}
As atividades desta etapa envolveram três temas de aprendizagem: a organização do código em um repositório, a divisão do trabalho em tarefas e prazos e a preparação de uma base de imagens para uso posterior. A separação dos conjuntos também introduziu a distinção entre dados destinados ao treinamento e dados reservados para verificação. No contexto educativo, foi necessário definir uma linguagem acessível para os materiais.

O registro disponível descreve as atividades, mas não apresenta uma reflexão individual sobre os conceitos compreendidos, as dúvidas ou os problemas enfrentados.

Permanecem pendentes os links de acesso ao código e ao quadro, a inclusão do painel de amostras e as reflexões individuais.

\subsubsection{Evidências da Sprint 1}
Os artefatos abaixo foram mencionados no registro de execução. Seus links e imagens devem ser incorporados para permitir a consulta do trabalho realizado.
\begin{itemize}
\item \textbf{Repositório:} \pendente{inserir o endereço do GitHub e a referência da versão correspondente ao encerramento da sprint.}
\item \textbf{Preparação dos dados:} script \path{preparacao_dados.py}. \pendente{Inserir o link para o arquivo e, se disponível, o registro de sua execução.}
\item \textbf{Painel de amostras:} imagem exportada no diretório \path{artefatos/}. \pendente{Inserir o painel e seu nome completo.}
\item \textbf{Jira:} \pendente{inserir uma captura ou exportação do quadro com as sprints, os prazos, os responsáveis e a situação das tarefas no encerramento.}
\end{itemize}

\subsubsection{Encaminhamento para a próxima sprint}
A continuidade prevista é estudar uma implementação introdutória de GAN, adaptar o código e executar um primeiro treinamento. Esses encaminhamentos seguem o planejamento entregue; não representam atividades concluídas na Sprint 1.

\subsection{Segunda Sprint: Primeira GAN e Alinhamento Educativo}
\label{sec:sprint2}
\textbf{Período:} 09/10 a 30/10/2026. \textbf{Peso da entrega:} 5\%.

O objetivo desta etapa foi estudar e adaptar uma implementação de GAN, executar um primeiro treinamento e registrar exemplos visuais. O registro técnico da Sprint 2~\cite{registroS2} documenta o desenvolvimento e as verificações descritos a seguir. A entrega técnica inicial foi alcançada.

\subsubsection{Implementação inicial e funcionamento da GAN}
O ponto de partida foi o commit \texttt{0523e6c}, de autoria do Allan, que contém a arquitetura, o código de treinamento e grades visuais identificadas como épocas 0, 1, 5 e 15. Essas imagens foram preservadas com o padrão de nomes \path{artefatos/sprint2_epoca_*.png}. O arquivo de progresso do treinamento e o log dessa execução não estavam disponíveis na cópia do repositório utilizada para a adaptação local. Por isso, não são atribuídos a essa execução valores de duração ou perdas.

O gerador recebe números aleatórios e os transforma em imagens. O discriminador recebe imagens da base e imagens geradas, aprendendo a distinguir suas origens. As duas redes são atualizadas alternadamente: o sinal de erro obtido na avaliação das imagens orienta as atualizações do gerador. A implementação atual não recebe uma descrição textual nem permite escolher a categoria a ser gerada.

Para acompanhar o processo, a equipe passou a registrar os resultados por época, isto é, a cada passagem pelo conjunto de treinamento. As imagens são processadas em lotes. As perdas orientam o ajuste das redes, mas não representam porcentagens de acerto ou medidas diretas de qualidade visual.

\subsubsection{Adaptações e correção do uso dos dados}
\label{sec:ajustesS2}
O código inicial utilizava as 60.000 imagens do conjunto oficial de treinamento, sem aplicar a divisão de 50.000 para treinamento e 10.000 para validação preparada na Sprint 1. Essa diferença foi corrigida centralizando a divisão estratificada no arquivo \path{dados.py}, com semente 42. A conferência encontrou 5.000 imagens de cada categoria no treinamento e 1.000 na validação, sem sobreposição entre os subconjuntos.

O teste oficial, composto por outras 10.000 imagens, permaneceu fora do treinamento nas duas versões. Na execução local adaptada, tanto a validação quanto o teste ficaram fora da atualização das redes. As imagens de validação continuam reservadas, sem métrica quantitativa implementada nesta etapa.

A arquitetura original de \path{modelos.py} foi mantida. As principais adaptações foram:
\begin{itemize}
\item Configuração do número de épocas e do tamanho do lote pelo terminal.
\item Salvamento, a cada época, das duas redes, dos otimizadores e dos estados aleatórios necessários à retomada.
\item Exportação separada do gerador para produzir imagens sem repetir o treinamento.
\item Registro de parâmetros, ambiente, duração, perdas médias e imagens por época.
\item Visualização em modo de avaliação, com escala de cinza fixa entre 0 e 1, sem alterar as estatísticas de normalização da rede.
\item Separação dos arquivos da implementação inicial e da nova execução local.
\end{itemize}

O uso do mesmo ruído de entrada ao longo das épocas permite observar mudanças nas saídas para uma mesma entrada. As imagens antigas, entretanto, foram produzidas com ajuste automático de contraste por imagem e com a rede em modo de treinamento. Assim, as grades antigas e novas não constituem uma comparação controlada, pois também diferem no conjunto de treinamento e no procedimento de visualização.

\subsubsection{Execução local e verificações de funcionamento}
Em 30/10/2026, foi executado o comando:
\begin{quote}
\small\verb|python3 treinar_gan.py --epocas 1 --nome sprint2_local|
\end{quote}
Foi concluída uma época completa sobre 50.000 imagens, em CPU, com lote de até 128 imagens e duas threads. Os valores medidos são apresentados na Seção~\ref{sec:resultadosS2}. O planejamento não estabelecia um número mínimo de épocas; essa execução atende ao objetivo de realizar um primeiro treinamento, sem encerrar a melhoria do modelo.

Após o treinamento, o gerador salvo foi carregado em outro processo pelo comando \texttt{python3 gerar\_imagens.py}. Essa execução exportou uma nova grade em \path{artefatos/geracao_local/imagens_geradas.png}, demonstrando a possibilidade de usar o modelo salvo sem treinar novamente.

Além da execução completa de uma época, foram realizados testes curtos com 256 imagens reais. O registro técnico informa:
\begin{itemize}
\item Execução da preparação dos dados e conferência da separação e do balanceamento dos conjuntos.
\item Salvamento e retomada do treinamento. Na comparação entre duas épocas contínuas e uma época seguida de retomada, os pesos das duas redes e as perdas foram idênticos no ambiente CPU testado.
\item Verificação de que a visualização não alterava pesos ou estatísticas da rede.
\item Carregamento do gerador salvo e exportação de novas imagens em outro processo.
\item Recusa de sobrescrita de uma execução existente, de retomada sem arquivo de progresso, de configurações incompatíveis e de números de épocas inválidos.
\end{itemize}

Esses testes verificam o funcionamento do código no ambiente utilizado. As execuções identificadas como \path{verificacao_*} não são evidências da qualidade visual de um treinamento completo, e a equivalência observada na retomada não foi generalizada para outros ambientes.

\subsubsection{Recursos, organização da equipe e aprendizados}
Foram utilizados Python, PyTorch, os módulos de preparação e treinamento do projeto, o repositório Git/GitHub e o Jira. A execução local registrada utilizou CPU com duas threads. O modelo do processador e a quantidade de memória não foram transcritos no resumo recebido; a configuração completa da execução está indicada no arquivo \path{configuracao.json}.

O Allan contribuiu com a arquitetura, o treinamento inicial e as grades preservadas no commit \texttt{0523e6c}. O Guilherme informou a inclusão das tarefas no backlog do Jira e o encerramento da Sprint 2. As adaptações locais e os testes foram documentados no registro técnico, mas sua distribuição individual e as horas de cada integrante ainda precisam ser consolidadas. Essa participação complementa as responsabilidades previstas inicialmente, sem restringir o estudo e o desenvolvimento a um único integrante.

Os principais conceitos trabalhados foram a interação entre gerador e discriminador, a diferença entre época e lote, a interpretação das perdas e a utilidade de salvar o modelo. Um ponto de atenção identificado foi que preparar uma divisão de dados não garante que o treinamento a utilize: foi necessário conectar essa divisão ao código de execução. Outro foi a influência da visualização na interpretação das imagens, o que motivou a padronização da escala e do modo de avaliação.

Esses pontos descrevem o aprendizado associado às atividades da equipe. As reflexões individuais e as dificuldades de cada integrante ainda não foram registradas separadamente.

\subsubsection{Comparação entre o planejado e o realizado}
\begin{center}
\small
\begin{tabularx}{\textwidth}{@{}>{\raggedright\arraybackslash}p{0.29\textwidth}Y@{}}
\toprule
\textbf{Previsão da Sprint 2} & \textbf{Situação no encerramento} \\
\midrule
Estudar e adaptar uma GAN & Implementação inicial disponível; adaptação local manteve a arquitetura e corrigiu o uso dos dados. Conceitos básicos registrados. \\
Executar o primeiro treinamento & Uma época completa sobre 50.000 imagens concluída localmente, com parâmetros, tempo e perdas registrados. \\
Salvar e analisar imagens & Grades locais e nova geração com o modelo salvo registradas; análise visual descritiva disponível. Imagens originais preservadas separadamente. \\
Organizar tarefas e documentar & Tarefas inseridas no backlog do Jira e registro técnico incorporado ao relatório. Horas, reflexões individuais e links ainda precisam ser consolidados. \\
\bottomrule
\end{tabularx}
\end{center}

A Sprint 2 foi encerrada com a primeira implementação executável e resultados iniciais de geração. O encerramento do período não significa que todas as atividades previstas tenham sido atendidas: as pendências de documentação seguem para acompanhamento.

\subsubsection{Evidências e encaminhamentos}
O registro técnico indica os seguintes arquivos da execução local:
\begin{itemize}
\item \path{artefatos/sprint2_local/configuracao.json}: parâmetros e ambiente.
\item \path{artefatos/sprint2_local/treinamento.log}: progresso e duração.
\item \path{artefatos/sprint2_local/historico.csv}: épocas concluídas e perdas médias.
\item \path{artefatos/sprint2_local/epoca_*.png}: imagens geradas com ruído fixo.
\item \path{artefatos/sprint2_local/perdas.png}: gráfico das perdas de treinamento.
\item \path{checkpoints/sprint2_local/ultimo.pt}: estado para retomada, mantido fora do Git.
\item \path{checkpoints/sprint2_local/gerador.pt}: gerador exportado, mantido fora do Git.
\end{itemize}

Os caminhos acima identificam os artefatos relatados; os arquivos de imagem e os links do repositório e do Jira ainda não foram incorporados a este documento. A análise visual apresentada na seção seguinte corresponde ao relato da execução, sem reprodução das grades nesta versão.

Para a Sprint 3, os encaminhamentos são continuar o treinamento, organizar a evolução visual, iniciar a interface e o roteiro educativo e preparar a apostila. A retomada da execução local até 15 épocas está preparada pelo comando:
\begin{quote}
\small\verb|python3 treinar_gan.py --epocas 15 --nome sprint2_local --retomar|
\end{quote}
Nesse caso, o total pretendido é de 15 épocas, com mais 14 a executar após a época já concluída. Esse prolongamento não foi realizado no registro de encerramento da Sprint 2. A interface educativa e a apostila permanecem para desenvolvimento posterior.

\section{Resultados Obtidos e Verificação dos Materiais}
\subsection{Resultados Parciais da Sprint 1}
Os resultados registrados na Sprint 1 são o repositório organizado, o quadro inicial do Jira, o script de preparação do Fashion-MNIST, a divisão dos dados e o painel de exemplos da base. Também foram descritos o objetivo educativo e os leitores pretendidos. As horas informadas somam 5 horas de trabalho ativo. Permanecem pendentes a incorporação das evidências e as reflexões individuais.

\subsection{Resultados do Treinamento Local da Sprint 2}
\label{sec:resultadosS2}
A execução local de 30/10/2026 produziu o primeiro registro completo de treinamento disponível para esta atualização, reunindo configuração, perdas e duração no relato técnico~\cite{registroS2}.

\begin{center}
\small
\begin{tabularx}{\textwidth}{@{}>{\raggedright\arraybackslash}p{0.47\textwidth}Y@{}}
\toprule
\textbf{Item} & \textbf{Resultado informado} \\
\midrule
Identificação da execução & \path{sprint2_local} \\
Imagens utilizadas no treinamento & 50.000 \\
Épocas completas & 1 \\
Tamanho máximo do lote & 128 imagens \\
Processamento & CPU, com duas threads \\
Tempo do laço de treinamento & 296,8 segundos (4 min 56,8 s) \\
Perda média do discriminador & 0,7726 \\
Perda média do gerador & 1,5397 \\
Uso do modelo salvo & Gerador carregado em outro processo e nova grade exportada \\
\bottomrule
\end{tabularx}
\end{center}

O tempo medido não inclui o download dos dados nem a exportação dos artefatos. As perdas são registros do treinamento e não representam acurácia, porcentagem de imagens corretas ou uma comparação direta de desempenho entre as duas redes.

\subsection{Evolução Visual e Limitações}
Segundo a análise descrita no registro técnico, as saídas da época 0 eram texturas acinzentadas, sem formas reconhecíveis. Após uma época, apareceram contornos que lembram roupas e calçados, mas ainda com borramento, padrões quadriculados e imagens difíceis de interpretar. O resultado indica uma mudança inicial das saídas, ainda com limitações de nitidez.

As grades do trabalho inicial do Allan, identificadas como épocas 0, 1, 5 e 15, são evidências de outra execução. Não devem ser associadas ao tempo ou às perdas da execução local. Também não se deve apresentar a imagem antiga de 15 épocas como evolução posterior à época local, pois os procedimentos de treinamento e visualização diferem.

A partir dos 296,8 segundos medidos na primeira época, uma extrapolação simples para mais 14 épocas resulta em aproximadamente 69 minutos adicionais de processamento. Essa estimativa está sujeita à carga e à temperatura da máquina e não representa uma medição de treinamento já realizado. A retomada foi verificada apenas nos testes curtos descritos anteriormente.

\subsection{Situação da Proposta Educativa}
Até o encerramento da Sprint 2, foram obtidos o código executável, o gerador salvo e imagens iniciais que poderão apoiar as próximas etapas. A interface educativa e a apostila ainda não foram implementadas nesta etapa.

A conclusão desta sprint é, portanto, um avanço técnico inicial. A continuidade deverá transformar esse resultado em uma demonstração acessível, mantendo as limitações visíveis e organizando os aprendizados da equipe.

\begin{thebibliography}{9}
\addcontentsline{toc}{section}{Referências}

\bibitem{fashionmnist}
XIAO, Han; RASUL, Kashif; VOLLGRAF, Roland.
\textbf{Fashion-MNIST: a Novel Image Dataset for Benchmarking Machine Learning Algorithms}.
2017. Repositório do projeto disponível em:
\url{https://github.com/zalandoresearch/fashion-mnist}.

\bibitem{registroS2}
EQUIPE DO PROJETO.
\textbf{Sprint 2 --- implementação e primeiro treinamento da GAN}.
Registro técnico de desenvolvimento, com resultado da execução de 30 out. 2026.
Documento interno fornecido para atualização do relatório.
\end{thebibliography}

\end{document}